\section{Screening and affine energies at a fixed stage}
\label{sec:screening}

Throughout this section the starting level $n$ is fixed and sufficiently
large. Constants bearing a subscript $n$ may depend arbitrarily on $n$,
as well as on the fixed $h,E$. They do not depend on the torus size once
that size is sufficiently large. We use the stage-$n$ construction of
\Cref{hier:sampling}. We also allow the prescribed artificial probes:
at each prescribed site there is at most one such probe, of strength
$\tau$ or $\tau\vartheta(v_x)$, where $0\leq\tau\leq1$.
A mask specifies which of these choices is made at each site.
Masks, observation grids, and test centers may depend on the hierarchy
offset and on external independent randomness, but not on the potentials
or on the branch decisions. When conditioning on the offset below, we
also fix any such external randomness. Uniform assertions about these
prescriptions mean estimates for every individual prescription, not a
single event on which all prescriptions work simultaneously.

Write $\Lambda_r(x)$ for the max-norm ball of radius $r$, with shortest
coordinate distances on the torus. A tile of level $l>n$ is
\emph{hazardous} if its selector is present, regardless of acceptance,
or if one of its own corner labels fails the level-$l$ validity test.
The construction gives, conditional on the offset,
\begin{equation}
 \PP(U\text{ is hazardous}\mid\text{offset})
 \leq C b_l^{-K_0},\qquad K_0=10^6.
 \label{screen:hazard}
\end{equation}
Consequently, for a square of radius $A$, the probability that a
hazardous tile of side greater than $D$ meets that square is at most
\begin{equation}
 C\sum_{b_l>D}(1+A/b_l)^2b_l^{-K_0}.
 \label{screen:hazard-count}
\end{equation}
All the counting estimates below are used within a small fixed fraction
of the torus side, where squares may be unwrapped in the usual way.

\subsection{Prospective observations and a generic observation bound}

Choose a large power of three $Q_1$, and tile the torus by observation
cores of side $Q_1$. Set $\beta=1/12$ and require $Q_1^\beta\geq b_n$.
For a core centered at $z$, test candidate level-$n$ roots in
$\Lambda_{3Q_1}(z)$ using the prospective rule from
\Cref{hier:sampling}: retain all local acceptance rules, but omit
ancestor selectors and validity tests at sides exceeding $Q_1^\beta$.
The resulting test for one candidate uses randomness within distance
$C Q_1^\beta$ of that candidate, conditional on the offset.

Each test succeeds with probability at least $c q_n p_n>0$.
To see this, first expose the required selector statuses and the needed
corner-box identities. Conditional on the candidate being eligible and
all required labels being valid, its acceptance likelihood integrates
to one, so its selector and acceptance contribute $q_np_n$.
The probability of the required strict-ancestor selector absences and
the own and ancestor validity conditions is bounded below by a positive
absolute constant: their failures have a summable union bound along the
ancestor chain, made small by the choice of the starting level.
This also explains why the local likelihood is not additionally
conditioned on finer corner data.

Cover $\Lambda_{2Q_1}(z)$ by mesh squares of side comparable to
$Q_1^{1/4}$. In each such square choose candidates separated by a
sufficiently large multiple of $Q_1^\beta$. There are at least
$c Q_1^{2(1/4-\beta)}=cQ_1^{1/3}$ candidates with disjoint randomness
neighborhoods, for sufficiently large $Q_1$ at fixed $n$.
Their success indicators are independent conditional on the offset.
The probability that some mesh square contains no successful root is
therefore at most
\begin{equation}
 C Q_1^{3/2}\exp(-c_n Q_1^{1/3}).
 \label{screen:net-failure}
\end{equation}
In particular, with arbitrarily high probability, the prospective
observation roots form a $C Q_1^{1/4}$-net of
$\Lambda_{2Q_1}(z)$, also for nearest-neighbor graph distance after
changing the absolute constant.

The next argument applies to every possible observation set with this
net property. Its equations always use the true potentials.

\begin{lemma}[Generic local observation]
\label{screen:observation}
For sufficiently large $Q_1$, almost every choice of the true
potentials on $\Lambda_{3Q_1}(z)$ has the following property.
For every set $\mathcal O\subset\Lambda_{3Q_1}(z)$ that is a
$C Q_1^{1/4}$-net of $\Lambda_{2Q_1}(z)$, there is a finite constant
$C_{\mathrm{obs}}$ such that
\begin{equation}
 \norm{u|_{\mathrm{core}}}
 \leq C_{\mathrm{obs}}
 \left(\norm{(H-E)u|_{\Lambda_{3Q_1}(z)}}
             +\norm{u|_{\mathcal O}}\right)
 \label{screen:observation-bound}
\end{equation}
for every vector on $\Lambda_{3Q_1+1}(z)$.
After $Q_1$ is fixed, a deterministic $C_{\mathrm{obs}}$ can be chosen
so that this bound holds for every such set $\mathcal O$ outside an
event of arbitrarily small probability. The exceptional probability
is uniform in the offset and in all sufficiently large tori.
\end{lemma}

\begin{proof}
It suffices to show that a vector annihilated by both constraints on the
right of \Cref{screen:observation-bound} vanishes on the core.
All these constraints are real, so a counterexample would have a real
counterexample. Fix its support pattern
$\mathcal P\subset\Lambda_{3Q_1+1}(z)$ and an observation-set pattern
$\mathcal O$, and suppose that $\mathcal P$ meets the core.
Set $\ell=\lceil C'Q_1^{1/4}\rceil$, with $C'$ large enough for the
graph-distance net property.

For any fixed sufficiently small $c_0>0$, some integer
$r\in[Q_1,2Q_1]$ satisfies
\begin{equation}
 \abs{\mathcal P\cap
       (\Lambda_{r+1}(z)\setminus\Lambda_{r-\ell}(z))}
 \leq c_0\ell^{-2}
       \abs{\mathcal P\cap\Lambda_{r-\ell}(z)}.
 \label{screen:thin-support-annulus}
\end{equation}
Indeed, if this failed at every radius, apply its reverse inequality at
radii separated by $\ell+1$. At each step the support count would grow
by a factor greater than $1+c_0\ell^{-2}$. The first inner square
contains the nonempty core support, and there are at least
$cQ_1/\ell$ steps. The resulting count would be at least
$\exp(cQ_1/\ell^3)$, whereas the available neighborhood has only
$O(Q_1^2)$ sites. Since $Q_1/\ell^3$ grows as a positive multiple of
$Q_1^{1/4}$, this is impossible for large $Q_1$.

For a radius satisfying \Cref{screen:thin-support-annulus}, write
\[
 d_0'=\abs{\mathcal P\cap\Lambda_{r-\ell}(z)},\qquad
 p=\abs{\mathcal P\cap\Lambda_r(z)},\qquad
 m=\abs{\mathcal P\cap\Lambda_{r+1}(z)}.
\]
Every deep support site $x\in\mathcal P\cap\Lambda_{r-\ell}(z)$ has
an observed zero within graph distance $\ell$. Follow a shortest path
to such a zero and take its first vertex outside $\mathcal P$.
This is a missing site in $\Lambda_r(z)$ adjacent to the support.
At that site the equation $(H-E)u=0$ is a nonzero homogeneous linear
equation in the $m$ coordinates on
$\mathcal P\cap\Lambda_{r+1}(z)$; its diagonal term vanishes.

Select a maximal collection of these missing sites with pairwise graph
distance greater than two. Their neighbor sets are disjoint, and each
row is nonzero, so the corresponding linear equations are independent.
By maximality, each deep site is within distance $\ell+2$ of a selected
missing site. A graph ball of this radius has at most $C\ell^2$ sites.
Thus the number $k$ of independent equations satisfies
\[
 k\geq c\ell^{-2}d_0'.
\]
On the other hand,
$m-p\leq c_0\ell^{-2}d_0'$ by
\Cref{screen:thin-support-annulus}. Choosing $c_0<c$ gives
\begin{equation}
 m-k<p.
 \label{screen:dimension-count}
\end{equation}

The equations at the $p$ nonzero sites in $\Lambda_r(z)$ determine
their true potentials through
\[
 v_x=E-\frac{\sum_{y:\,|y-x|_1=1}u_y}{u_x}.
\]
The admissible nonzero coordinates lie in a linear space of dimension
at most $m-k$. On the set where the relevant denominators are nonzero,
these formulas define a smooth map into $\R^p$. Cover its domain by
countably many bounded patches on which the denominators are bounded
away from zero. On each patch the map is Lipschitz, and its image is
Lebesgue null by \Cref{screen:dimension-count}. For completeness, a
bounded patch in dimension $a<p$ may be covered by $O(\varepsilon^{-a})$
cubes of side $\varepsilon$; their Lipschitz images have total
$p$-dimensional covering volume $O(\varepsilon^{p-a})$, which tends
to zero.
The joint law of the $p$ true potentials has a density, so this support
and observation pattern has probability zero. There are only finitely
many patterns and admissible radii. Taking their union proves the
kernel assertion, including for an observation set chosen from the
potentials.

Let $A$ be the map consisting of the residual and observation
constraints, and let $B$ be restriction to the core. The assertion
$\ker A\subset\ker B$ makes $B$ a well-defined linear map on
$\operatorname{im}A$ through $A$. Finite dimensionality gives a finite
operator bound, which is \Cref{screen:observation-bound}; complex
vectors obey the same bound by complexification.
The maximum of these finite bounds over all possible net sets is
finite almost surely. Taking a sufficiently high quantile gives a
deterministic $C_{\mathrm{obs}}$. The local true product law is
unchanged by offset or by passage to a larger torus, proving the last
assertion. No deterministic unique-continuation assertion is used.
\end{proof}

\subsection{Projected resolvents and shield comparisons}

The decay estimate below uses an exponential-conjugation argument of
Combes--Thomas type \cite{CombesThomas1973}. We first establish the
projected inverse bound on the observed coordinates; no bound on the
full inverse is required.

Call an observation core \emph{good} if the prospective roots form
the specified net, the deterministic observation bound holds, and no
hazardous tile of side greater than $Q_1^\beta$ meets
$\Lambda_{4Q_1}$ about its center. The prospective roots of a good
core are actual stage-$n$ ports: all omitted ancestor requirements
hold. Moreover, every real reference change in its tested rows occurs
at a cell port of side at most $Q_1^\beta$. Such a port has strength
bounded below by a positive number depending on $Q_1$.

We may count only a subset of these good cores. In a partly thinned
or Dirichlet-restricted system, we count only cores for which the
physical neighborhood $\Lambda_{3Q_1+1}(z)$ remains present, and for
which all cell ports, their strengths, and the real reference values
in the tested rows $\Lambda_{3Q_1}(z)$ agree with the full stage-$n$
system. We call these tests \emph{intact}. The neighbor layer preserves
every hopping term in the true equations; artificial probes may still
have any of the allowed masks.
Let $J_{\mathrm g}$ be the coordinate embedding of the sites in the
counted cores, and set
\[
 \Pi=[J_{\mathrm g}\ \ Y],\qquad
 G=(H^\#-E-iYY^\top)^{-1},
\]
where $Y$ includes every current port, whether or not it lies in a
counted core.

\begin{proposition}[Projected screening]
\label{screen:projected}
For all the systems just described,
\begin{equation}
 \norm{\Pi^\top G\Pi}\leq C_n.
 \label{screen:projected-bound}
\end{equation}
Give cost one to nearest-neighbor edges whose endpoints both belong
to counted good cores, and cost zero to other edges; any of the
positive costs may also be dropped. If two source and response sectors
of $\Pi$ have cost distance at least $t$, their resolvent block has
norm at most $C_n\exp(-c_nt)$.

Suppose, further, that two systems agree on an inner region and a
surrounding shield, including all ports there, and that counted good
tests in the shield are usable in both systems. Count costs only on
shared good-good edges in the shield. If every crossing of the shield
from the inner region to the exterior costs at least $t$, then the
responses between common inner ports differ by at most
$C_n\exp(-c_nt)$ in operator norm. This includes comparison with a
Dirichlet restriction buffered beyond the shield.
\end{proposition}

\begin{proof}
Put $u=G\Pi w$, and abbreviate
$A=\norm{J_{\mathrm g}^\top u}$,
$B=\norm{Y^\top u}$, and $W=\norm w$.
In each tested neighborhood the true residual is
\[
 (H-E)u=\Pi w+(H-H^\#)u+iYY^\top u.
\]
There are at most two ports at a site and their strengths are uniformly
bounded, so $\norm Y$ is uniformly bounded. The real changes are
bounded and, within these rows, supported at cell ports with the
positive strength lower bound just noted. Thus their contribution is
bounded by $C_nB$. The observed values obey the same bound, using the
accepted level-$n$ ports. The testing neighborhoods have bounded
overlap. Summing the squared observation inequalities gives
\begin{equation}
 A\leq C_n(W+B).
 \label{screen:residual-observation}
\end{equation}
Taking the imaginary part of the inverse equation, or using the Ward
identity of \Cref{mat:ward}, yields
\[
 B^2\leq (A+B)W\leq C_n(W+B)W.
\]
This implies $B\leq C_nW$, and then $A\leq C_nW$ by
\Cref{screen:residual-observation}. It proves
\Cref{screen:projected-bound} without a bound on the full inverse.

Let $F$ be a real function on physical sites with
$|F(x)-F(y)|$ at most the assigned cost of every nearest-neighbor
edge. Conjugate the site matrix by multiplication by $e^{sF}$.
Diagonal entries do not change, and every changed hopping entry has
both endpoints among the counted good sites. Hence
\[
 e^{sF}(H^\#-E-iYY^\top)e^{-sF}
 =H^\#-E-iYY^\top+J_{\mathrm g}D_sJ_{\mathrm g}^\top,
 \qquad \norm{D_s}\leq C(e^s-1).
\]
For a sufficiently small fixed $s>0$, depending on $n$, we have
$\norm{D_sJ_{\mathrm g}^\top GJ_{\mathrm g}}<1/2$.
The inverse difference formula and
\Cref{screen:projected-bound} therefore bound the conjugated projected
inverse by another $C_n$. Assign to each column of $\Pi$ the value of
$F$ at its physical site, so that repeated columns at a site have the
same weight. If $F_\Pi$ is this diagonal matrix, the conclusion is
\begin{equation}
 \norm{e^{sF_\Pi}\Pi^\top G\Pi e^{-sF_\Pi}}\leq C_n.
 \label{screen:weighted-bound}
\end{equation}
Choose $F$ to be cost distance from the source sector. Restriction of
\Cref{screen:weighted-bound} to a response sector at distance $t$
gives the stated exponential attenuation.

For the two-system comparison, use this same distance from the inner
region, with costs counted only in the shield. For $t$ larger than an
absolute constant, choose a multiplier $\chi$ equal to one where
$F\leq t/3$, zero where $F\geq2t/3$, and affine in between.
It is supported in the common region, and changes only across counted
good-good edges. Cut off the first system's columns sourced at common
inner ports by $\chi$. In the second inverse equation the intended
source is unchanged, and the only error is a hopping commutator
supported on good coordinates. Each edge of that commutator has
$F\geq t/3-1$ at both endpoints. The weighted bound in the first
system makes the error, as a good-coordinate source, at most
$C_ne^{-c_nt}$ in operator norm. Correcting it with the second
system's projected inverse costs at most another $C_n$.
Taking the common inner-port response proves the comparison. Bounded
$t$ is covered by the unweighted projected bounds after enlarging the
constant. The same calculation applies to a buffered Dirichlet
restriction, because the cutoff vanishes before its boundary.
\end{proof}

\subsection{Crossing costs and the wall variant}

By \Cref{screen:net-failure,screen:observation}, the probability of a
local test failure can be made arbitrarily small, first by taking
$Q_1$ large and then $C_{\mathrm{obs}}$ large. The additional hazard
cost for one core is at most
\[
 C\sum_{b_l>Q_1^\beta}(1+Q_1/b_l)^2b_l^{-K_0},
\]
which tends to zero as $Q_1$ grows. Consequently, in any neighborhood
of radius a fixed multiple of $Q_1$, the probability that some core
meeting that neighborhood is bad is arbitrarily small. All these
probabilities are conditional on the offset, uniformly in that offset.

\begin{proposition}[Full-stage path estimate]
\label{screen:paths}
There is an absolute $\alpha>0$ such that, for the full stage-$n$
good-core costs and all sufficiently large $r$ up to a small fixed
fraction of the torus side,
\begin{equation}
 \PP\left(\begin{array}{c}
 \text{some nearest-neighbor crossing from }\Lambda_r(x)
       \text{ to }\Lambda_{2r}(x)^c\\
 \text{has cost less than }c_nr^\alpha
 \end{array}\,\middle|\,\text{offset}\right)
 \leq C_nr^{-50}.
 \label{screen:path-bound}
\end{equation}
Only costs of edges with both endpoints in the intervening band need
be counted. The constants are uniform in the center and the allowed
prescriptions.
\end{proposition}

\begin{proof}
Fix a sufficiently large absolute integer $B_1$, and use the integer
radii
\[
 r_0=100Q_1,\qquad r_j=B_1^jr_0,
\]
as long as $r_j\leq b_N/100$. Let $p_j'$ be the worst conditional
probability, over offsets, centers, and sufficiently large tori, of a
crossing of the $r_j$-band with cost less than $2^j$.
Trim a crossing at its two radial boundaries, so that the tested
segment lies entirely in the band.

A crossing of the $r_{j+1}$-band with cost less than $2^{j+1}$
forces two failures at scale $r_j$ at well-separated centers from a
bounded mesh list. Indeed, take successive visits to radii rounded
from
$5r_{j+1}/4$, $3r_{j+1}/2$, and $7r_{j+1}/4$, and replace each
visited site by a mesh center within distance $r_j$.
For $B_1$ large enough, the corresponding $r_j$-annuli lie in the
larger band, are pairwise disjoint, and are crossed by subsegments of
the original path. Their centers are more than $20r_j$ apart.
If two of these crossings each had cost at least $2^j$, the original
cost would be at least $2^{j+1}$. Thus at least two are failures.
The number of possible mesh-center pairs is an absolute constant
after $B_1$ has been fixed.

Temporarily truncate the hazard tests used for these smaller annuli
at side $r_j$. Each resulting failure event then depends only on
randomness within distance at most $5r_j$ of its center. This includes
the local prospective tests, the true observation equations, and
the complete corner triples and selectors of every tested tile.
The two separated events are independent conditional on the offset.
Truncation can only enlarge the good set, so each truncated failure
probability is at most $p_j'$.
If no omitted hazardous tile meets the needed neighborhoods, the
truncated and full costs coincide. The union bound
\Cref{screen:hazard-count} therefore gives
\begin{equation}
 p_{j+1}'\leq C(p_j')^2+
 C\sum_{b_l>r_j}(1+r_{j+1}/b_l)^2b_l^{-K_0}
 \leq C(p_j')^2+C r_j^{-K_0}.
 \label{screen:path-recursion}
\end{equation}
The constants here may depend on the fixed absolute $B_1$.

Choose a small absolute $a>0$ with
$Ca\leq B_1^{-50}/2$. The base tests above allow $p_0'\leq a$.
Increase $Q_1$, if necessary, so that the hazard term at the base is
at most $aB_1^{-50}/2$. Since $K_0>50$, induction in
\Cref{screen:path-recursion} gives
\[
 p_j'\leq aB_1^{-50j}.
\]
Set $\alpha=\log_{B_1}2$; thus
$2^j=(r_j/r_0)^\alpha$. This proves the assertion at recursion radii.
For a general sufficiently large radius $r$, take the largest $r_j$
with $r_j\leq r/100$. A crossing of the $r$-band has a translated
$r_j$-band crossing wholly inside it, at one of a bounded mesh list
of centers. Here $r_j$ is comparable to $r$, with absolute comparison
constants. The union bound gives
\Cref{screen:path-bound}, after changing $c_n,C_n$.
This construction also covers noninteger radii and uses no costs at
the outer band boundaries.
\end{proof}

\begin{proposition}[Screening on retained walls]
\label{screen:walls}
Fix an aligned square grid of side $L_*$, deterministic given the
offset. Suppose a system retains the full stage-$n$ data on a wall
set containing strips of width $\delta L_*$ along every grid boundary,
where $\delta>0$ is a fixed absolute constant. Count only good cores
whose centers are at least $5Q_1$ inside this wall set, so that all
their tests are intact. There is an absolute constant $A_{\mathrm w}$
such that, for $L_*$ sufficiently large at fixed $n$ and
\[
 A_{\mathrm w}L_*\leq r\leq c b_N,
\]
every crossing of the $r$-band has wall cost at least $c_n'r^\alpha$
outside an event of conditional probability at most $C_n'r^{-30}$.
Here $c_n',C_n'$ and the threshold on $L_*$ may depend on $n$;
$A_{\mathrm w}$ depends only on the absolute strip fraction and
geometric constants. In particular, the constants in the bound are
uniform as $L_*$ grows. As before, only edges within the crossed band
need be counted.
\end{proposition}

\begin{proof}
It suffices to count good cores whose centers are at least $5Q_1$
inside the prescribed boundary strips. These cores are among those
allowed in the statement, and dropping the other costs can only
decrease crossing costs. This counted geometric region is deterministic
conditional on the offset.

Let $r_0^{\mathrm w}=K L_*$ for a sufficiently large fixed absolute
$K$. A crossing of its band contains a smaller annulus crossing
deep inside one wall strip. To verify this, unwrap coordinates about
the large-band center $x$ and stop the path when it first reaches
radius approximately $7r_0^{\mathrm w}/4$. One of its signed
coordinates attains this value. Between signed coordinate distances
$5r_0^{\mathrm w}/4$ and $3r_0^{\mathrm w}/2$ from $x$, choose a
parallel grid boundary line. Such a line exists because $K$ is large.
The path must visit this line, while its other coordinate has modulus
at most $7r_0^{\mathrm w}/4+1$.

Put $\rho=\lfloor c_{\mathrm w}L_*\rfloor$ for a small fixed
$c_{\mathrm w}>0$. Choose a mesh center within $\rho/10$ of this
visit. The entire $\rho$-annulus about that center lies both inside
the original large band and inside the wall strip; choosing
$c_{\mathrm w}$ sufficiently small and then $L_*$ sufficiently large
also leaves the $5Q_1$ buffer for all relevant core tests.
By its later visit to radius $7r_0^{\mathrm w}/4$, the path has left
the small annulus. It therefore provides the claimed crossing.
Since $r_0^{\mathrm w}/\rho$ is bounded by an absolute constant,
only boundedly many mesh centers are needed. Applying
\Cref{screen:paths} at those centers gives a base wall threshold
$a_nL_*^\alpha$, with failure probability at most $C_nL_*^{-50}$.

Run the same recursion as in \Cref{screen:path-recursion}, now with
the deterministic wall restriction at every scale. The three-annulus
geometry is unchanged. After truncating hazards, separated events
still use disjoint randomness neighborhoods, since the counted strips
add no randomness conditional on the offset. At radius
$r_j^{\mathrm w}=B_1^jK L_*$ the threshold is
\[
 2^j a_nL_*^\alpha
   =a_nK^{-\alpha}(r_j^{\mathrm w})^\alpha.
\]
The failure recurrence is again a square of the preceding probability
plus $O((r_j^{\mathrm w})^{-K_0})$.
For a sufficiently large lower threshold on $L_*$, its base bound
$C_nL_*^{-50}$ and this recurrence imply
$C_n'(r_j^{\mathrm w})^{-30}$ at all recursion radii: at a new radius
$B_1r$, the two terms are $O_n(r^{-60})$ and $O(r^{-K_0})$, which
are bounded by $C_n'(B_1r)^{-30}$ once the base radius is large.
This induction uses constants independent of $L_*$.
Finally use a comparable smaller recursion scale and a bounded mesh
list inside a general band, as in \Cref{screen:paths}. Taking
$r\geq100r_0^{\mathrm w}$ suffices. This is an absolute multiple of
$L_*$, as claimed.
\end{proof}

\subsection{Jump tails and fixed-stage limits}

At full stage $n$, with the possible extra probes, define
\begin{equation}
 d_{n,N}(x)=\frac12
 \sum_{a:\,\operatorname{site}(a)=x}\sum_j
       |S_{ja}|^2\dist_2(x,\operatorname{site}(j))^2.
 \label{screen:local-energy}
\end{equation}
Here $\dist_2$ is Euclidean distance formed from the shortest
coordinate displacements. Let $d_{n,N}^{\leq D}(x)$ and
$d_{n,N}^{>D}(x)$ denote the same sum restricted respectively to jump
lengths at most $D$ and greater than $D$.

\begin{proposition}[Fixed-stage limits]
\label{screen:limits}
For each fixed sufficiently large $n$ there is a finite number
$d_n\geq0$ with the following properties.
The jump tails satisfy
\begin{equation}
 \sup\E\left[d_{n,N}^{>D}(x)\mid\text{offset}\right]
       \leq\omega_n(D),\qquad \omega_n(D)\longrightarrow0,
 \label{screen:jump-tail}
\end{equation}
uniformly in the site, sufficiently large tori, prescribed masks, and
$0\leq\tau\leq1$.
For a deterministic site $x$, with expectation including the uniform
offset,
\begin{equation}
 \E d_{n,N}(x)\longrightarrow d_n
       \qquad (N\longrightarrow\infty,\ \tau\longrightarrow0),
 \label{screen:point-limit}
\end{equation}
uniformly in the prescriptions.
If $A$ is a growing square, whose position may be chosen from the
offset and external independent randomness, then
\begin{equation}
 \frac1{|A|}\sum_{x\in A}d_{n,N}(x)\longrightarrow d_n
       \quad\text{in mean and in probability}.
 \label{screen:spatial-limit}
\end{equation}
The estimates for this convergence are uniform conditional on the
offset, over such square positions and prescriptions, as the square
side and the torus size grow and $\tau\to0$.
Bounded-length jumps capture these spatial averages to any given
tolerance with arbitrarily high probability if the jump cutoff is
chosen first and sufficiently large.
\end{proposition}

\begin{proof}
At most two ports occur at a site: one accepted cell root and one
artificial probe. Thus the source rank in
\Cref{screen:local-energy} is at most two. For large $D$, use a
separating annulus about $x$ with radius comparable to $D/200$.
Its cost is at least $c_nD^\alpha$ outside an event of conditional
probability at most $C_nD^{-50}$, by \Cref{screen:paths}.
The projected attenuation in \Cref{screen:projected} applies from
the source ports at $x$ to all ports at distance greater than $D$
simultaneously. Unitarity bounds that block on the exceptional event.
The source-rank bound gives
\begin{equation}
 \E\left[\sum_{a:\,\operatorname{site}(a)=x}
     \sum_{j:\,\dist_2(x,\operatorname{site}(j))>D}|S_{ja}|^2
        \,\middle|\,\text{offset}\right]
 \leq C_n\bigl(e^{-c_nD^\alpha}+D^{-50}\bigr).
 \label{screen:mass-tail}
\end{equation}
Multiply the bound at dyadic distances $2^qD$ by
$C(2^qD)^2$ and sum. The series is convergent and tends to zero
with $D$, proving \Cref{screen:jump-tail}.
This also controls jumps comparable to the torus side: their
separating annulus can still be taken at radius $D/200$, within the
allowed small torus fraction. There is no final unestimated shell.

We next construct local approximants, keeping $n$ fixed. Fix a jump
cutoff $D$. First restrict the effective site data to a Dirichlet
square of a sufficiently large fixed radius $R_{\mathrm{loc}}$
about $x$. A shield buffered inside this square encloses $x$ and all
targets at distance at most $D$. The shield comparison of
\Cref{screen:projected} and the path estimate give arbitrarily small
expected error in the truncated sum as $R_{\mathrm{loc}}$ grows.
On its exceptional event use unitarity; the truncated sum is bounded
by $CD^2$. The comparison concerns a block with source rank at most
two, so its error does not acquire a factor from the number of remote
ports.

Next choose a fixed superscale $b_k$ much larger than
$R_{\mathrm{loc}}$ and omit all hierarchy rules above that scale
when forming the restricted effective data. Use the prospective
ancestor convention, the same true potentials, and the same
selectors and acceptance uniforms below the cutoff.
The restricted effective data agree with the original data unless an
omitted hazardous tile meets the restriction. Indeed, absent such a
hazard, no omitted selector can suppress a retained root, and every
omitted ancestor validity gate is satisfied. The local likelihoods
do not otherwise use ancestor-label identities. The probability of
disagreement is bounded by
\[
 C\sum_{b_l>b_k}(1+R_{\mathrm{loc}}/b_l)^2b_l^{-K_0},
\]
which tends to zero with $k$. A retained reference value and its
likelihood use only their finite tile. The resulting truncated
Dirichlet quantity therefore depends on randomness within distance
$O(R_{\mathrm{loc}}+b_k)$ of $x$.

For this fixed finite approximant we may remove artificial probes
by sending $\tau$ to zero. To keep its channel space fixed, pad it
with one artificial channel at every site of the restriction and
give absent channels strength zero. Such a column is zero in $Y$
and contributes only an isolated identity entry to $S$.
Its inverse at $\tau=0$ exists almost
surely: a kernel vector of the passive matrix vanishes at cell-port
sites, where all real reference changes are supported, and hence
would be a kernel vector of the true Dirichlet restriction minus
$E$. The latter has nonzero determinant almost surely by the joint
diagonal density. Inverse continuity now gives
$S\to S_{\mathrm{cell}}\oplus1$ on this fixed channel space.
In particular, the cell-port scattering block converges.
Mixed entries involving an artificial
probe tend to zero, as do off-diagonal entries between artificial
probes. Their diagonal entries do not contribute to
\Cref{screen:local-energy}, since their jump length is zero.
Bounded convergence applies to the truncated sums, which are
uniformly bounded by $CD^2$.
There are only finitely many local masks, including the choices of
constant or profile strength, and finitely many offset patterns
modulo the fixed superscale. Taking a maximum over these finitely
many choices makes this convergence uniform. This step does not
require any uniform lower bound for $\vartheta(v_x)$.

Combining these constructions with the jump-tail bound proves the
following useful precise statement. For every $\varepsilon>0$ there
is a bounded no-extra local function $Z_x$, with fixed dependence
radius $\rho_{\mathrm{loc}}=O(R_{\mathrm{loc}}+b_k)$, such that
\begin{equation}
 \E[|d_{n,N}(x)-Z_x|\mid\text{offset}]<\varepsilon
 \label{screen:local-approximation}
\end{equation}
for every sufficiently large $N$ and sufficiently small $\tau$,
uniformly in masks, offsets, and sites. The choices of $D$,
$R_{\mathrm{loc}}$, and $k$ are made in that order, before these
limits.

Conditional on the offset, the $Z_x$ have finite dependence range:
two collections farther apart than $2\rho_{\mathrm{loc}}$ are
independent. Their conditional means are periodic at the superscale
$b_k$, and the pattern for a different offset is a translate of the
same pattern. Thus their period-averaged mean, say $\mu_Z$, is
independent of offset. It is also their mean at any deterministic
site after averaging over the uniform offset. For all sufficiently
large tori these local laws are identical. By
\Cref{screen:local-approximation}, the numbers
$\E d_{n,N}(x)$ are Cauchy in the joint limit, uniformly in the
prescriptions. Their finite limit is $d_n$, and
$|\mu_Z-d_n|\leq\varepsilon$.

For a square $A$ of side $a$, covering its interior by complete
$b_k$-periods gives
\[
 \left|\frac1{|A|}\sum_{x\in A}
       \E[Z_x\mid\text{offset}]-\mu_Z\right|
 \leq C\norm{Z}_\infty\frac{b_k}{a}.
\]
Counting pairs with nonzero conditional covariance also gives
\[
 \operatorname{Var}\left(\frac1{|A|}\sum_{x\in A}Z_x
                  \,\middle|\,\text{offset}\right)
 \leq C\norm{Z}_\infty^2\frac{\rho_{\mathrm{loc}}^2}{a^2}.
\]
Both bounds hold for every square position, including one chosen
from the offset. The same bounded-range counting applies on the
torus, with periodic distances. Average
\Cref{screen:local-approximation} over $A$, take the square and
torus sizes large, and then send $\varepsilon$ to zero. This proves
the asserted conditional uniform convergence in mean and hence in
probability. Notice that a conditional point mean need not equal
$d_n$; it is the period average that removes its dependence on the
offset pattern.
Finally, the expected difference between the full and truncated
spatial averages is at most $\omega_n(D)$, even conditional on the
offset. Markov's inequality, followed by
\Cref{screen:spatial-limit}, proves the bounded-jump assertion.
\end{proof}

\subsection{Compact profiles and convergence of affine energies}

\begin{proposition}[Compact-profile estimate]
\label{screen:profiles}
Let $\varphi:\R^2\to\R^q$ be compactly supported and Lipschitz,
piecewise smooth with bounded gradients and piecewise regular
interfaces. Let its center $c$ be prescribed from the offset and
external independent randomness. At full stage $n$,
\begin{equation}
 \limsup_{\substack{R\to\infty,\ \tau\to0\\b_N/R\to\infty}}
 \E\,\mathcal Q_{\mathrm{start}}
       \bigl(\varphi((\cdot-c)/R)\bigr)
 \leq C d_n\int_{\R^2}|\nabla\varphi|^2\,\dd x,
 \label{screen:profile-bound}
\end{equation}
where $C$ is absolute. Here $n$ and the profile are fixed; the
thresholds may depend arbitrarily on them. The estimate is uniform
over the allowed masks and centers, and remains valid with the
expectation conditioned on the offset.

More generally, let $f$ be supported in a square of side at most
$A R$ whose position is prescribed from the offset and external
independent randomness, and suppose that $f$ has global Lipschitz
constant at most $B/R$. Then
\begin{equation}
 \E\,\mathcal Q_{\mathrm{start}}(f)
       \leq C_{A,B}(d_n+o(1))
 \label{screen:rough-profile}
\end{equation}
in the same fixed-$n$ regime, uniformly over all such data and
prescriptions. Thus each assertion holds to any fixed additive
tolerance for all sufficiently large spatial and torus scales and
all sufficiently small $\tau$; no rate uniform in $n$ is asserted.
\end{proposition}

\begin{proof}
Fix a large jump cutoff $D$. If a profile difference is nonzero,
at least one endpoint belongs to its support. By symmetry of $S$,
the contribution of jumps longer than $D$ is bounded by
\[
 \frac{C\operatorname{Lip}(\varphi)^2}{R^2}
       \sum_{x\in\operatorname{supp}\varphi_R}
                    d_{n,N}^{>D}(x),
 \qquad \varphi_R(x)=\varphi((x-c)/R).
\]
There are $O(R^2)$ sites in this support. Its expected value tends
uniformly to zero as $D\to\infty$, by
\Cref{screen:jump-tail}.

For jumps of length at most $D$, the squared profile difference is
bounded by $R^{-2}$ times the squared jump length times a local
supremum of $|\nabla\varphi|^2$. Away from the interfaces, these
suprema converge to $|\nabla\varphi|^2$ in the Riemann-sum limit.
The sites whose rescaled $D$-neighborhood meets an interface have
vanishing relative area, and the global Lipschitz bound controls
their contribution. Thus the short-jump energy is bounded by
\[
 \frac{C}{R^2}\sum_x w_R(x)d_{n,N}^{\leq D}(x),
\]
where $w_R$ is uniformly bounded, supported in an $O(R)$ square,
and has Riemann-sum limit $\int|\nabla\varphi|^2$.

To justify averaging these weights also conditional on the offset,
use the fixed local approximants from
\Cref{screen:local-approximation}. Their bounded periodic conditional
means have period average arbitrarily close to $d_n$.
On complete periods, the slowly varying Riemann-sum weights may be
replaced by a representative value with error tending to zero;
the incomplete boundary periods contribute vanishingly. The
uniform conditional expected approximation error contributes at
most $C\varepsilon R^{-2}\sum_x w_R(x)$.
This proves the bound for the short jumps after sending the spatial
and torus scales to infinity and $\tau$ to zero. Then let the local
approximation error and the jump-tail error tend to zero. All
finite cutoffs and local neighborhoods were chosen at fixed $n$
before taking these limits.

For the more general data, use the global Lipschitz bound directly.
Symmetry charges every nonzero difference to the supporting square,
so
\[
 \mathcal Q_{\mathrm{start}}(f)
 \leq \frac{CB^2}{R^2}
      \sum_{x\text{ in the supporting square}}d_{n,N}(x).
\]
The uniform spatial mean estimate in \Cref{screen:limits} now gives
\Cref{screen:rough-profile}. This argument only uses support size
and the Lipschitz bound, so it is uniform for varying profiles with
the stated bounds.
\end{proof}

To approximate affine coordinate data on a torus without choosing a
discontinuous coordinate branch, put
\begin{equation}
 f_N(x)=\frac{b_N}{2\pi}
       \bigl(\cos(2\pi x_i/b_N),\sin(2\pi x_i/b_N)\bigr)_{i=1,2}
       \in\R^4.
 \label{screen:affine-data}
\end{equation}
Its chord distances satisfy
\begin{equation}
 \frac2\pi\dist_2(x,y)\leq|f_N(x)-f_N(y)|\leq\dist_2(x,y),
 \label{screen:chord-comparison}
\end{equation}
and their ratio to $\dist_2(x,y)$ tends to one uniformly for bounded
nonzero jumps as $N\to\infty$. Indeed, for shortest coordinate
displacements $\delta_i$, the squared chord distance is
$\sum_i(b_N\sin(\pi\delta_i/b_N)/\pi)^2$.
Truncate the jumps, use this local limit, and then use
\Cref{screen:jump-tail} for the tails. Translation after offset
averaging and \Cref{screen:point-limit} give
\begin{equation}
 b_N^{-2}\E\,\mathcal Q_{S_{n,N}}(f_N)\longrightarrow d_n.
 \label{screen:affine-limit}
\end{equation}
There are no artificial probes in this identity.

\begin{proposition}[Limit of the affine energy densities]
\label{screen:limit}
The fixed-stage limits satisfy
\begin{equation}
 d_n\longrightarrow d\quad\text{for some }0\leq d<\infty.
 \label{screen:d-limit}
\end{equation}
More precisely, a deterministic $\varepsilon_n\to0$ can be chosen
such that
\begin{equation}
 d_m\leq(1+\varepsilon_n)d_n+\varepsilon_n
             \qquad(m>n)
 \label{screen:almost-monotone}
\end{equation}
for all sufficiently large $n$, uniformly in $m$.
\end{proposition}

\begin{proof}
Apply \Cref{hier:comparison} without extras to $f_N$ in the full
transition from stage $n$ to stage $m$. The data have a uniformly
bounded Lipschitz constant by \Cref{screen:chord-comparison}.
The area-normalized charge estimate of \Cref{hier:charges} has a
vanishing $n$-tail independent of $m,N$, as does the multiplicative
error in the comparison. Divide by $b_N^2$, fix $n<m$, and send
$N$ to infinity using \Cref{screen:affine-limit} for both stages.
This proves \Cref{screen:almost-monotone} with errors uniform in
$m>n$; no limit uniform in the pair $(n,m)$ was taken.

Fixing one sufficiently large $n_0$ in that inequality first bounds
$d_m$ for all $m>n_0$. Let $d$ be the finite nonnegative lower limit
of this bounded sequence and choose $n_k\to\infty$ with
$d_{n_k}\to d$. For each $k$, \Cref{screen:almost-monotone} gives
\[
 \limsup_{m\to\infty}d_m
    \leq(1+\varepsilon_{n_k})d_{n_k}+\varepsilon_{n_k}.
\]
Sending $k\to\infty$ proves that the upper limit is at most $d$.
This proves convergence.
\end{proof}

The order of parameters in this section will be used explicitly
below. One fixes $n$ before choosing observation scales, quantiles,
jump cutoffs, or local approximation radii. Their possibly very
large values impose only fixed thresholds on subsequent spatial
scales. With fixed additional parameters, one can then take $m>n$
so large that scales of order $\exp(b_m^2)$ exceed every such
threshold. A probe strength $\tau=C_0/R$ at those scales tends to
zero and eventually satisfies
$\tau\leq\exp(-c_*b_m)$, as required by
\Cref{hier:comparison}. Finally $N$ is increased after these choices.
Neither this procedure nor \Cref{screen:d-limit} supplies or uses a
rate of convergence uniform in $n$.
